\expected{bikram-v-evolation}

\item \textbf{Rule hierarchies}. Like many areas of law, copyright law operates
based on a set of interrelated rules that, ideally, can be arranged in a
hierarchical form. Upon encountering any legal rule, it is always worth making
sure you understand where the rule fits into a larger doctrinal hierarchy. Cases
should lay out the full rule hierarchy, but sometimes as in the \emph{Bikram
Yoga} case, the court forgets to do this. So here's a brief overview.

\begin{figure}

\begin{flowchart}
\blocks
node (infr) {Copyright\\infringement}

node[above right=of infr] (ownership) {Ownership of\\valid copyright}
node[above right=of ownership] (102a) {\S~102(a) subject\\matter}
node[right=of ownership] (ownership-other) {Other rules\ldots}

node[right=of infr] (copy) {Copying}

node[below right=of infr] (original) {Original\\elements}
node[above right=of original]  (independent) {Independent\\creation}
node[right=of original] (creative) {Minimal\\creativity}
node[below right=of original] (102b) {Exceptions:\\\S~102(b)}
node[right=of 102b] (idea) {Idea/expression\\dichotomy}
;
\conn[] infr.east to ownership.west;
\conn[] infr.east to copy.west;
\conn[] infr.east to original.west;
\conn[] ownership.east to 102a.west;
\conn[] ownership.east to ownership-other.west;
\conn[] original.east to independent.west;
\conn[] original.east to creative.west;
\conn[] original.east to 102b.west;
\conn[] 102b to idea.west;
\end{flowchart}

\caption{Tree representation of the copyright doctrinal hierarchy. The phrases
in the boxes are shorthand for the exact language of the elements of the rules,
as stated in statutes and case law. Due to space concerns, the ``top'' of the
tree is toward the left.}
\label{f:copyright-tree}
\end{figure}

The top of the hierarchy should always be the rule that determines the legal
relationship between the parties. For copyright law, that would be the rule for
copyright infringement: ``To establish infringement, two elements must be
proven: (1) ownership of a valid copyright, and (2) copying of constituent
elements of the work that are original.'' \sentence{feist-v-rural at 361}.
Observe that the second element of the test actually requires two distinct
things: (2a) copying, and (2b) original elements.

As the case illustrates, \S~102(a) lays out some of the requirements for valid
copyright ownership (there are others). No definition of copying is given; other
cases do that. \Inline{feist-v-rural} gives the following definition of original
elements: ``independently created by the author'' and ``possesses at least some
minimal degree of creativity.'' \sentence{feist-v-rural at 345}. But \S~102(b)
statutorily limits which ``original elements'' are relevant for copying, and
\emph{Bikram Yoga} describes \S~102(b) as embodying the ``idea/expression
dichotomy'' doctrine.

All these doctrines may be connected into a hierarchy, starting from the
infringement rule. \unskip\footnote{Note that this rule hierarchy is based on
the one given in \inline{feist-v-rural} and \emph{Bikram Yoga}, but it is not
the definitive copyright rule hierarchy. Among other things, some courts will
place the \S~102(b) and the idea/expression dichotomy as sub-elements of
copying, and there is controversy over whether the idea/expression dichotomy is
the whole test for \S~102(b) as opposed to just one of several tests. Law school
classes often portray rules as having fixed, well-defined elements, which is the
right place to start when learning how to work with legal rules, but be aware
that the real world is often messier.} A graphical ``tree'' form (the editor's
preference) is given in Figure~\ref{f:copyright-tree}, though it can also be
represented as an indented outline. The key is that rules are connected to their
elements, and then those elements are further connected to sub-rules and
sub-elements, and so on.

\item \textbf{Using doctrinal hierarchies}. A doctrinal hierarchy is used in two
main ways. First, when performing a legal analysis on an exam or in a brief or
memorandum, you typically should start from the top of the hierarchy and work
downwards, stating the rules as you proceed. This ensures that you touch on all
the relevant legal questions and gives you an order for addressing the issues.
So, for example, in \emph{Bikram Yoga}, the initial legal issue is whether
Evolation infringed Choudhury's copyright. That depends on whether Choudhury
owned a valid copyright, whether there was copying, and whether the copying was
of original elements. After ownership and copying have been addressed, the
presence of original elements is a sub-issue that depends on \S~102(b), which in
turn depends on the idea/expression dichotomy doctrine.

Second, once you reach a conclusion as to an item in the hierarchy, you work
your way upwards to determine the consequence. In \emph{Bikram Yoga}, the court
finds that the sequence of yoga poses is an idea rather than expression. What
consequence does that have for the evaluation of \S~102(b)? What consequence
does that have, in turn, for original elements, and then for infringement
overall? The doctrinal hierarchy is the roadmap for how any given legal
determination connects to the overall consequences for the parties, allowing you
to predict outcomes and advise clients.

\item Why should ideas be outside the scope of copyright protection? This is a
great question for testing some of the theories of property law identified in
\mref{intro-property-rationales}. After all, it often takes valuable labor to
come up with good ideas like a beneficial sequence of yoga poses (a Lockean
rationale), and society benefits from granting incentives for people to come up
with those ideas (an economic rationale). The Ninth Circuit offers quite a few
rebuttals, and perhaps you can come up with more. Consider, for each of them,
whether there is a property theory or intuition at work.

\defcase{kelley-v-chicago}{
Kelley v. Chicago Park District, 635 F.3d 290 (7th Cir. 2011)
}
\defcase{burrow-giles-v-sarony}{
Burrow-Giles Lithographic Co. v. Sarony, 111 U.S. 53 (1884)
}

\item What else belongs in or out of the scope of copyright protection? You will
learn the precise answers to these questions in a copyright law course, but
consider the following based on the reasoning of \emph{Bikram Yoga}, your
intuitions, and general principles of law:
\begin{itemize}
\HiNoId
\item A telephone book, alphabetically ordered. \sentence{see feist-v-rural}.
\item A beautifully manicured garden. \sentence{see kelley-v-chicago}.
\item A photograph, considering that it merely depicts the appearance of the
real world. \sentence{see burrow-giles-v-sarony}.
\item A selfie photograph taken by a monkey, or a painting-like image made with
generative artificial intelligence. \emph{See} note~\ref{n:ai-creator} in this
book.
\item A state's official publication of its statutory code. \emph{See}
\mref{georgia-v-public-resource}, later in this book.
\end{itemize}

\expecting{georgia-v-public-resource}
\expecting{erickson-v-trinity-theatre-qs}
